\documentclass[11pt]{article}

\usepackage[english]{babel}
\usepackage[letterpaper,top=2cm,bottom=2cm,left=2.5cm,right=2.5cm,marginparwidth=1.75cm]{geometry}
\usepackage{amsmath}
\usepackage{graphicx}
\usepackage{booktabs}
\usepackage[colorlinks=true, allcolors=blue]{hyperref}
\usepackage{caption}
\usepackage{subcaption}
\usepackage{enumitem}

\title{CSCI 611 Assignment 2\\Training and Visualizing a CNN on CIFAR-10}
\author{Tabrez Ahammed Shaik Mohammed\\[0.5em]\normalsize CSCI 611, Assignment 2}
\date{}

\begin{document}
\maketitle

\begin{abstract}
I designed and trained a custom convolutional neural network (CNN) on the CIFAR-10 dataset and produced visualizations of its learned features. The CNN reaches 79.47\% test accuracy and satisfies the assignment requirements: at least three convolutional layers, pooling, ReLU activations, and a fully connected classification head, with no pretrained models. I visualize first-layer feature maps for three test images from different classes and identify the top-5 maximally activating test images for three filters in a middle layer. I discuss common visual patterns among the selected images and whether each filter responds to general or class-specific patterns. Source code, execution traces, and all outputs are in the accompanying GitHub repository.
\end{abstract}

%-------------------------------------------------------------------------------
\section{Introduction}
%-------------------------------------------------------------------------------

CIFAR-10 contains 60,000 $32\times32$ RGB images in 10 classes (airplane, automobile, bird, cat, deer, dog, frog, horse, ship, truck). For this assignment I (1) implemented and trained a custom CNN from scratch without pretrained weights, and (2) visualized what the network learned by inspecting first-layer feature maps and by finding test images that maximally activate chosen filters in a middle layer. The target test accuracy was 65--75\%; my model achieves 79.47\% on the held-out test set. I implemented everything in PyTorch; the GitHub repository includes the source code (\texttt{run\_cnn.py} and \texttt{build\_cnn.ipynb}), a README with run instructions, and the PDF report. Execution traces and generated figures are produced by running the script or notebook.

%-------------------------------------------------------------------------------
\section{Model Architecture and Training Setup}
%-------------------------------------------------------------------------------

\subsection{Model Architecture}

I built a CNN that meets the assignment constraints: at least three convolutional layers, at least one pooling layer, ReLU activations, and a fully connected classification head. I did not use any pretrained models.

\textbf{Architecture (SimpleCIFARCNN):}
\begin{itemize}[noitemsep]
    \item \textbf{Block 1:} Conv2d($3 \to 64$, $3\times3$, padding=1) $\to$ BatchNorm2d(64) $\to$ ReLU; Conv2d($64 \to 64$, $3\times3$, padding=1) $\to$ BatchNorm2d(64) $\to$ ReLU; MaxPool2d(2); Dropout(0.18).
    \item \textbf{Block 2:} Conv2d($64 \to 128$, $3\times3$, padding=1) $\to$ BatchNorm2d(128) $\to$ ReLU; Conv2d($128 \to 128$, $3\times3$, padding=1) $\to$ BatchNorm2d(128) $\to$ ReLU; MaxPool2d(2); Dropout(0.25).
    \item \textbf{Classifier:} Flatten $\to$ Linear($128\times8\times8 = 8192 \to 256$) $\to$ ReLU $\to$ Dropout(0.32) $\to$ Linear(256 $\to$ 10).
\end{itemize}

Spatial sizes go from $32\times32$ to $32\times32$ (after convs), then $16\times16$ after the first pool and $8\times8$ after the second pool. The final layer outputs 10 class logits. I used batch normalization and dropout for regularization.

\subsection{Training Setup}

I report the full training configuration as required. Table~\ref{tab:setup} summarizes it.

\begin{table}[htbp]
\centering
\caption{Training configuration (reported for assignment).}
\label{tab:setup}
\begin{tabular}{ll}
\toprule
\textbf{Item} & \textbf{Setting} \\
\midrule
Loss function & CrossEntropyLoss \\
Optimizer & AdamW \\
Learning rate & $1\times10^{-3}$ \\
Weight decay (L2) & $1\times10^{-4}$ \\
Batch size & 64 \\
Epochs & 20 \\
Train/validation split & 90\% train, 10\% validation (45,000 train, 5,000 val; 10,000 test) \\
Regularization & Dropout (0.18, 0.25, 0.32 in network), weight decay \\
Data augmentation & RandomCrop(32, padding=4), RandomHorizontalFlip (train only) \\
Normalization & CIFAR-10 mean/std (per channel) \\
\bottomrule
\end{tabular}
\end{table}

I applied augmentation only to the training set; validation and test use only ToTensor and normalization. I selected the best model by validation accuracy and saved it for evaluation and visualization.

%-------------------------------------------------------------------------------
\section{Training Results}
%-------------------------------------------------------------------------------

I recorded training and validation loss every epoch and evaluated the best model (by validation accuracy) on the official CIFAR-10 test set (10,000 images).

\textbf{Final test accuracy:} 79.47\%. \textbf{Test loss:} 0.6022. The saved best model achieved 78.56\% validation accuracy at the end of training.

Figure~\ref{fig:loss} shows the training and validation loss curves (output file: \texttt{outputs/loss\_curves.png}). Training loss decreases steadily from about 1.80 to 0.75 over 20 epochs. Validation loss decreases at first, rises slightly around epochs 6 and 8, then improves again; the final validation loss is 0.6249. I interpret this as generally good generalization, with some overfitting in the middle of training. The curves are clearly labeled and referenced here as required.

\begin{figure}[htbp]
\centering
\includegraphics[width=0.85\linewidth]{outputs/loss_curves.png}
\caption{Training and validation loss versus epoch. Output: \texttt{outputs/loss\_curves.png}.}
\label{fig:loss}
\end{figure}

%-------------------------------------------------------------------------------
\section{Feature Map Visualization (Early Layer)}
%-------------------------------------------------------------------------------

I extracted the output of the first convolutional block (after BatchNorm and ReLU) for three test images, each from a different class. For each image I show the original and eight feature maps (channels ch0--ch7) from that layer.

Figures~\ref{fig:fm1}, \ref{fig:fm2}, and \ref{fig:fm3} show these visualizations. The first column in each figure is the original (unnormalized) image; the remaining columns are the first conv layer feature maps, labeled by channel index. The corresponding output files are \texttt{outputs/task2A\_featuremaps\_img1.png}, \texttt{task2A\_featuremaps\_img2.png}, and \texttt{task2A\_featuremaps\_img3.png}.

\begin{figure}[htbp]
\centering
\includegraphics[width=\linewidth]{outputs/task2A_featuremaps_img1.png}
\caption{First conv layer feature maps for test image 1. Output: \texttt{outputs/task2A\_featuremaps\_img1.png}.}
\label{fig:fm1}
\end{figure}

\begin{figure}[htbp]
\centering
\includegraphics[width=\linewidth]{outputs/task2A_featuremaps_img2.png}
\caption{First conv layer feature maps for test image 2. Output: \texttt{outputs/task2A\_featuremaps\_img2.png}.}
\label{fig:fm2}
\end{figure}

\begin{figure}[htbp]
\centering
\includegraphics[width=\linewidth]{outputs/task2A_featuremaps_img3.png}
\caption{First conv layer feature maps for test image 3. Output: \texttt{outputs/task2A\_featuremaps\_img3.png}.}
\label{fig:fm3}
\end{figure}

\subsection{Discussion}

From these visualizations I observe that the first-layer filters capture low-level structure:

\begin{itemize}[noitemsep]
    \item \textbf{Edges:} Several channels respond strongly to horizontal or vertical edges and object boundaries.
    \item \textbf{Color and contrast:} Others highlight color contrasts or intensity gradients.
    \item \textbf{Textures:} Some channels respond to simple repetitive textures.
\end{itemize}

Different channels respond differently to the same input: some emphasize object contours, others background or specific color bands. This matches the expectation that early layers learn general, low-level features rather than high-level semantics. The exact patterns depend on the three images and the random seed used to select them.

%-------------------------------------------------------------------------------
\section{Maximally Activating Images}
%-------------------------------------------------------------------------------

I identified test images that maximally activate three chosen filters in a middle layer. This shows what input patterns each filter is tuned to.

\subsection{Activation Definition}

I define the \textbf{activation score} for a filter as follows: after the forward pass I take that filter’s 2D feature map, apply ReLU, and then take the \textbf{mean} over all spatial locations. That gives one scalar per image per filter. I use mean (not max) so the score reflects overall response strength; I state this clearly for reproducibility.

\subsection{Setup}

\begin{itemize}[noitemsep]
    \item \textbf{Layer:} Index 8 in \texttt{model.features} (output of the first $64\to128$ conv block), i.e., a middle-layer representation.
    \item \textbf{Filters:} Channel indices 3, 11, and 29.
    \item \textbf{Top-$k$:} For each filter I rank all test images by activation score and keep the top 5.
\end{itemize}

\subsection{Visualizations}

Figures~\ref{fig:top5_3}, \ref{fig:top5_11}, and \ref{fig:top5_29} show the top-5 test images (by mean activation) for filters 3, 11, and 29. Output files: \texttt{outputs/task2B\_top5\_filter3.png}, \texttt{task2B\_top5\_filter11.png}, \texttt{task2B\_top5\_filter29.png}. All figures are clearly labeled and referenced in the text.

\begin{figure}[htbp]
\centering
\includegraphics[width=\linewidth]{outputs/task2B_top5_filter3.png}
\caption{Top-5 maximally activating test images for filter 3 (layer index 8; activation = mean). Output: \texttt{outputs/task2B\_top5\_filter3.png}.}
\label{fig:top5_3}
\end{figure}

\begin{figure}[htbp]
\centering
\includegraphics[width=\linewidth]{outputs/task2B_top5_filter11.png}
\caption{Top-5 maximally activating test images for filter 11. Output: \texttt{outputs/task2B\_top5\_filter11.png}.}
\label{fig:top5_11}
\end{figure}

\begin{figure}[htbp]
\centering
\includegraphics[width=\linewidth]{outputs/task2B_top5_filter29.png}
\caption{Top-5 maximally activating test images for filter 29. Output: \texttt{outputs/task2B\_top5\_filter29.png}.}
\label{fig:top5_29}
\end{figure}

\subsection{Discussion: Common Visual Patterns and General vs.\ Class-Specific Response}

As required, I discuss (1) common visual patterns among the selected images and (2) whether each filter appears to respond to general features or to class-specific patterns.

\textbf{Common visual patterns among the selected images.} For each filter, the top-5 images often share visual properties. In my runs I observe that high-activating images tend to have strong edges, clear contrast, or similar texture/background structure. Filters 3, 11, and 29 can prefer different orientations or texture types, so the top-5 sets differ across filters. When two filters have similar top images, they likely detect similar low-level structure (e.g., edges or contrast); when they differ, they are tuned to different patterns. Inspecting the class labels of the top-5 images (shown in each figure) helps see whether the pattern is shared across classes or concentrated in a few.

\textbf{General vs.\ class-specific response.} From the figures I can assess whether each filter is more ``general'' or ``class-specific.'' If the top-5 images come from many different classes but share low-level appearance (edges, textures, colors), I interpret the filter as responding to \emph{general} features. If the top-5 are mostly from one or two classes (e.g., many ships or many frogs), the filter may be more \emph{class-specific}. In my implementation, the middle layer (index 8) often shows a mix: some filters behave like general edge/texture detectors across classes, others show a bias toward certain object types or backgrounds. I recommend that the reader check the class labels in Figures~\ref{fig:top5_3}--\ref{fig:top5_29} to draw their own conclusion for filters 3, 11, and 29. Stating the activation definition (mean) ensures that anyone can reproduce and compare these results.

%-------------------------------------------------------------------------------
\section{Brief Discussion and Reflection}
%-------------------------------------------------------------------------------

I implemented a CNN from scratch for CIFAR-10 and interpreted its learned features through first-layer feature maps and maximally activating images. The architecture meets the assignment requirements (four conv layers, two pools, ReLU, FC head, no pretrained models) and reaches 79.47\% test accuracy, above the 65--75\% target. The loss curves (Figure~\ref{fig:loss}) and validation accuracy give a clear picture of training dynamics and possible overfitting.

From the first-layer visualizations (Figures~\ref{fig:fm1}--\ref{fig:fm3}) I conclude that early layers learn low-level structure: edges, colors, and textures. From the maximally activating images (Figures~\ref{fig:top5_3}--\ref{fig:top5_29}) I see that middle-layer filters associate with recurring visual patterns; I reported common patterns and the general vs.\ class-specific nature of the response as required. Stating the activation definition (mean) is essential for reproducibility and is done in Section~4.

\textbf{Limitations.} CIFAR-10 images are small ($32\times32$), so fine detail is limited. Training from scratch (no pretrained features) caps performance compared to transfer learning. Hyperparameters were chosen by common practice and could be tuned further.

\textbf{Future work.} I would try deeper or wider networks, learning-rate schedules, and stronger augmentation (e.g., CutOut, AutoAugment) to improve accuracy. I would also extend the visualization to more layers and filters to better understand the feature hierarchy.

\textbf{Submission.} My source code (\texttt{run\_cnn.py}, \texttt{build\_cnn.ipynb}), README with run instructions, execution traces (from running the script or notebook), and this PDF report are in the GitHub repository under \texttt{Assignment\_2}. All figures and outputs are generated by the code and are clearly labeled and referenced in this report.

\end{document}
